For \(d\in\mathcal D\), abbreviate the identified bridge contrast by
\[
 \theta_d(P):=\mu_{d,1}-\mu_{d,0},
\]
and let \(\mathcal F_m:=\sigma(\mathcal P_m)\).  Assumption \(\mathrm{ass{:}independent\text{-}pilot}\) makes \(\widehat d_m\), \(\widehat\alpha_m\), \(N_{E,B}\), and \(N_{O,B}\) measurable with respect to \(\mathcal F_m\) and independent of the main-study observations.  The pilot-oracle-regret theorem, the premise \(\underline v_m=o(1)\), and Assumption \(\mathrm{ass{:}nondegenerate\text{-}components}\) imply that there is an \(a_0\in(0,1/2)\) such that
\[
 \Pr\{a_0\le\widehat\alpha_m\le1-a_0\}\longrightarrow1.                \tag{1}
\]
Indeed, the finite set of population shares \(\alpha_d^*(P)\) is contained in the interior of \((0,1)\), and the pilot theorem gives uniform convergence of the estimated shares.  On the event in (1), both realized counts are of order \(B\).  Uniformly over \(d\), the allocation conclusion of that theorem gives
\[
 s_{0,B}^2-\mathcal V_{\widehat d_m}^*(P)=o_P(1),
 \qquad
 s_{0,B}^2:=B\left\{
 \frac{V_{E,\widehat d_m}}{N_{E,B}}
 +\frac{V_{O,\widehat d_m}}{N_{O,B}}
 \right\}.                                                               \tag{2}
\]

We next adapt the fixed-design limit to the pilot-selected observed bridge functional.  For fold \(\ell\), set
\[
 e_{h,d,\ell}:=\widehat h_{d,B}^{(-\ell)}-h_d,
 \qquad
 e_{q,d,\ell}:=\widehat q_{d,B}^{(-\ell)}-q_d.
\]
The exact doubly robust score algebra, using the outcome-bridge equation and the adjoint identity, leaves the conditional-mean remainder
\[
\begin{aligned}
 \mathcal R_{d,\ell}
 &=\sum_{a=0}^1(-1)^{1-a}
 \mathbb E_P\!\left[
 \frac{\mathbf 1\{A=a\}}{p_{O,a}}
 e_{q,d,\ell}\,(T_de_{h,d,\ell})\mathrel{\Big|}G=O\right]\\
 &=\sum_{a=0}^1(-1)^{1-a}
 \mathbb E_P\!\left[
 \frac{\mathbf 1\{A=a\}}{p_{O,a}}
 (T_d^\star e_{q,d,\ell})e_{h,d,\ell}\mathrel{\Big|}G=O\right].
\end{aligned}                                                              \tag{3}
\]
Treatment overlap, Cauchy--Schwarz, and the two representations in (3) give a finite constant \(C_P\), independent of \(d\), \(\ell\), and \(B\), such that
\[
 |\mathcal R_{d,\ell}|
 \le C_P\min\left\{
 \|e_{h,d,\ell}\|_{w,h,d}\|e_{q,d,\ell}\|_{2,O,q,d},
 \|e_{h,d,\ell}\|_{2,O,h,d}\|e_{q,d,\ell}\|_{w,q,d}
 \right\}.                                                               \tag{4}
\]
The constant is uniform because the design family and treatment set are finite.  Assumptions \(\mathrm{ass{:}uniform\text{-}outcome\text{-}weak\text{-}rate}\), \(\mathrm{ass{:}uniform\text{-}outcome\text{-}strong\text{-}rate}\), \(\mathrm{ass{:}uniform\text{-}selection\text{-}weak\text{-}rate}\), \(\mathrm{ass{:}uniform\text{-}selection\text{-}strong\text{-}rate}\), and \(\mathrm{ass{:}uniform\text{-}bridge\text{-}product\text{-}rate}\) make (4) \(o_P(B^{-1/2})\), uniformly over \(d\) and \(\ell\).  Conditional on each training fold, cross-fitting centers the remaining empirical nuisance terms.  The four vanishing-rate assumptions make their conditional \(L_2\) norms \(o_P(1)\), so their empirical averages are \(o_P(B^{-1/2})\).  This is the empirical-negligibility argument in \(\mathrm{lem{:}ikmw\text{-}fixed\text{-}design\text{-}limit}\), now uniform over the finite family.  Together with \(\mathrm{lem{:}designwise\text{-}gradient}\), it yields, uniformly over \(d\in\mathcal D\) and source allocations in the compact interval in (1),
\[
 \widehat\tau_{d,B}-\theta_d(P)
 =\frac1{n_E}\sum_{i=1}^{n_E}\phi_{E,d}(O_{E,i})
  +\frac1{n_O}\sum_{i=1}^{n_O}\phi_{O,d}(O_{O,i})
  +o_P(B^{-1/2}).                                                          \tag{5}
\]
Uniformity here needs no additional empirical-process premise: if it failed, finiteness of \(\mathcal D\) and compactness of the allocation interval would produce a subsequence with a fixed design and a convergent allocation, contradicting the fixed-design lemma.  Pilot independence therefore permits substitution of \(d=\widehat d_m\) and of the pilot-measurable source counts in (5).

Conditional on \(\mathcal F_m\), the two leading sums are independent and mean zero, and their variance after multiplication by \(B\) is \(s_{0,B}^2\).  The standardized fixed-design limit, the same finite-subsequence argument, and averaging over \(\mathcal F_m\) give
\[
 \frac{\sqrt B\{\widehat\tau_B-\theta_{\widehat d_m}(P)\}}
 {s_{0,B}}\rightsquigarrow\mathcal N(0,1).                               \tag{6}
\]

On the event in (1), Assumptions \(\mathrm{ass{:}main\text{-}gradient\text{-}variance\text{-}rate}\) and \(\mathrm{ass{:}main\text{-}variance\text{-}rate\text{-}vanishes}\) give
\[
\begin{aligned}
 |\widehat s_B^2-s_{0,B}^2|
 &\le \frac B{N_{E,B}}
 |\widetilde V_{E,\widehat d_m,B}-V_{E,\widehat d_m}|\\
 &\quad+\frac B{N_{O,B}}
 |\widetilde V_{O,\widehat d_m,B}-V_{O,\widehat d_m}|\\
 &=O_P(r_{\widetilde V,B})=o_P(1).
\end{aligned}                                                              \tag{7}
\]
Equations (2) and (7) prove both asserted variance conclusions.  Nondegeneracy and positive costs bound the finitely many values \(\mathcal V_d^*(P)\) away from zero, so \(\widehat s_B/s_{0,B}\to_P1\).  To transfer (6) to the estimated scale directly, fix a continuity point \(x\) of the standard normal distribution and \(0<\delta<1\).  On \(\{|\widehat s_B/s_{0,B}-1|\le\delta\}\), the distribution function after division by \(\widehat s_B\) is squeezed between the corresponding distribution functions after division by \(s_{0,B}\), evaluated at \(x(1-\delta)\) and \(x(1+\delta)\), with the inequalities reversed when \(x<0\).  The complement has probability tending to zero.  First taking \(m,B\to\infty\) in (6), and then \(\delta\downarrow0\), proves
\[
 \frac{\sqrt B\{\widehat\tau_B-(\mu_{\widehat d_m,1}-\mu_{\widehat d_m,0})\}}
 {\widehat s_B}\rightsquigarrow\mathcal N(0,1).                          \tag{8}
\]

By the displayed definition of \(I_B\), coverage of the pilot-selected bridge contrast is the event
\[
 \left\{
 \left|
 \frac{\sqrt B\{\widehat\tau_B-(\mu_{\widehat d_m,1}-\mu_{\widehat d_m,0})\}}
 {\widehat s_B}
 \right|\le z_{1-\eta/2}
 \right\}.
\]
Writing \(\Phi\) for the standard normal distribution function, convergence in (8) at the two continuity points \(\pm z_{1-\eta/2}\) gives
\[
 P\{\mu_{\widehat d_m,1}-\mu_{\widehat d_m,0}\in I_B\}
 \longrightarrow
 \Phi(z_{1-\eta/2})-\Phi(-z_{1-\eta/2})=1-\eta.
\]
Finally, if \(\mu_{d,1}-\mu_{d,0}=\tau\) for every \(d\in\mathcal D\), then the selected bridge contrast equals \(\tau\) for every pilot realization.  Substitution into (8) and the preceding coverage display proves the causal specialization.